\documentclass[11pt]{article}
\usepackage[T1]{fontenc}
\usepackage[utf8]{inputenc}
\usepackage{lmodern}
\usepackage{microtype}
\usepackage[margin=1in]{geometry}
\usepackage{amsmath,amssymb,amsthm,booktabs}
\usepackage{xurl}
\usepackage[colorlinks=true,linkcolor=blue,citecolor=blue,urlcolor=blue]{hyperref}

\theoremstyle{plain}
\newtheorem{theorem}{Theorem}[section]
\newtheorem{proposition}[theorem]{Proposition}
\newtheorem{lemma}[theorem]{Lemma}

\theoremstyle{remark}
\newtheorem{remark}[theorem]{Remark}

\DeclareMathOperator{\diag}{diag}
\DeclareMathOperator{\Res}{Res}

\title{Main-Eigenvalue Switching for Families of Diameter-Five Trees}
\author{Claude Make No Mistakes}
\date{November 2026}

\begin{document}
\maketitle

\begin{abstract}
We construct sign switchings that make every distinct adjacency eigenvalue main
for all uniform reflected diameter-five trees $T_r^p(a)$ ($r,a\geq 1$,
$p\geq 0$), allowing arbitrary numbers of direct leaves at both centers. We
also establish sufficient leaf-group size bounds for heterogeneous and
nonreflected diameter-five trees, and derive an exact system of finite
algebraic constraints for the remaining diameter-five cases. The proofs
separate leaf-difference and branch-contrast modes from the radial subspace and
resolve the remaining finite families by resultant congruences, modular
certificates, and exact arithmetic.
\end{abstract}

\section{Introduction}\label{sec:intro}

An eigenvalue $\lambda$ of the adjacency matrix $A$ of a finite simple graph
$G$ is called \emph{main} if its eigenspace is not orthogonal to the all-ones
vector $\mathbf{1}$. A \emph{sign switching} replaces $A$ by
$D_s A D_s$, where $D_s = \diag(s)$ for a sign vector
$s \in \{-1,1\}^{|V(G)|}$. Problem~1.6 of Akbari, Kumar, Mohar, and
Pragada~\cite{AKMP} asks whether every tree and every connected regular graph
of order at least three admits a sign switching for which every distinct
eigenvalue is main. Dong~\cite{Dong} established this for all trees of diameter
at most four, while Cha and Shan~\cite{ChaShan} studied diameter-five trees
and their uniform reflected families $T_r(a)$ and $T_r^p(a)$ in connection with
main eigenvalues.

Since $D_s^2 = I$, $D_s A D_s (D_s v) = D_s A v$, and
$\mathbf{1}^T D_s v = s^T v$, the matrices $A$ and $D_s A D_s$ share the same
distinct eigenvalues, and a distinct eigenvalue $\lambda$ of $D_s A D_s$ is main
if and only if $A$ has a $\lambda$-eigenvector $v$ with $s^T v \neq 0$. We say
that $s$ \emph{detects} $\lambda$ when such an eigenvector exists, and call $s$
a \emph{full switching} if $s$ detects every distinct eigenvalue of $A$.

In Section~\ref{sec:uniform} we construct a full switching for every uniform
reflected diameter-five tree $T_r(a)$ ($r,a \geq 1$) without direct leaves
(Theorem~\ref{thm:uniform}). In Section~\ref{sec:radial} we analyze the radial
invariant subspace of an arbitrary diameter-five tree and prove sufficient
leaf-group size bounds for both nonreflected trees
(Proposition~\ref{prop:heterogeneous}) and reflected trees with unequal branch
sizes (Proposition~\ref{prop:reflected}). In Section~\ref{sec:direct} we combine
the reflected bound with exact modular resultant certificates to prove the
existence of a full switching for every uniform reflected tree $T_r^p(a)$ with
direct leaves (Theorem~\ref{thm:uniform-direct}). Section~\ref{sec:remaining}
records the exact algebraic constraints governing arbitrary diameter-five trees,
and Section~\ref{sec:repro} summarizes the accompanying Python/SymPy and Lean~4
verification artifacts.

\section{Uniform reflected trees without direct leaves}\label{sec:uniform}

For integers $r,a \geq 1$, let $T_r(a)$ be the diameter-five tree obtained by
joining two centers by an edge and attaching $r$ branch vertices, each adjacent
to $a$ leaves, to each center.

\begin{theorem}\label{thm:uniform}
For every pair of integers $r,a\geq 1$, the tree $T_r(a)$ admits a sign
switching that makes every distinct adjacency eigenvalue main.
\end{theorem}

\subsection{Spectrum of \texorpdfstring{$T_r(a)$}{T\_r(a)}}\label{sec:spectrum-tra}

The tree $T_r(a)$ has $2 + 2r(a+1)$ vertices. We decompose its adjacency
spectrum into three mutually orthogonal invariant subspaces:

\begin{enumerate}
\item \emph{Leaf differences (eigenvalue $0$).} When $a \geq 2$, vectors
  supported on the $a$ leaves of a single branch whose coordinates sum to $0$
  are $0$-eigenvectors of $A$. Across all $2r$ branches they span a space of
  dimension $2r(a-1)$.

\item \emph{Branch contrasts (eigenvalues $\pm\sqrt{a}$).} When $r \geq 2$,
  choose weights $t_1,\dots,t_r$ on the $r$ branches of one side with
  $\sum_{i=1}^r t_i = 0$. For $x = \pm\sqrt{a}$, assign $x t_i$ to branch
  vertex $i$, $t_i$ to each of its $a$ leaves, and $0$ to both centers and the
  opposite side. Direct verification shows $A v = x v$, yielding a subspace of
  dimension $2(r-1)$ at each of $+\sqrt{a}$ and $-\sqrt{a}$.

\item \emph{Radial subspace (dimension $6$).} Vectors constant on each of the
  three vertex types (center, branch, leaf) on each side form a
  $6$-dimensional $A$-invariant subspace orthogonal to both spaces above. Split
  it into reflection parity $e \in \{+1,-1\}$, so that coordinates on the second
  side are $e$ times those on the first. Let $(c, b, \ell)$ denote the
  first-side center, branch, and leaf values of a parity-$e$ eigenvector with
  eigenvalue $x$. The leaf and branch equations read $b = x\ell$ and
  $c + a\ell = xb$. If $\ell = 0$, then $b = c = 0$, so every nonzero radial
  eigenvector has $\ell \neq 0$. Normalizing $\ell = 1$ gives
  $(c, b, \ell) = (x^2-a, x, 1)$, and the center equation becomes
  \begin{equation}\label{eq:fe}
    f_e(x) = (x-e)(x^2-a) - rx = 0.
  \end{equation}
\end{enumerate}

In the orthonormal basis of normalized type indicators, the restriction of $A$
to each parity-$e$ radial subspace is the symmetric unreduced tridiagonal
matrix with diagonal $(e,0,0)$ and off-diagonal entries $\sqrt{r},\sqrt{a} > 0$.
Its first-row recurrence determines every eigenvector from its first coordinate,
so each $f_e$ has three real, simple roots. The two cubics share no roots
because $f_+(x) - f_-(x) = -2(x^2-a)$ while
$f_e(\pm\sqrt{a}) = \mp r\sqrt{a} \neq 0$; the same evaluation shows that no
radial eigenvalue equals $\pm\sqrt{a}$. Moreover $f_e(0) = ea \neq 0$, so $0$
is not a radial eigenvalue. The dimensions
$2r(a-1) + 4(r-1) + 6 = 2 + 2r(a+1)$ exhaust the entire space, giving the
complete spectrum of $T_r(a)$.

\subsection{Switching construction and finite exceptions}\label{sec:switching-tra}

\begin{proof}[Proof of Theorem~\ref{thm:uniform}]
Start with every sign equal to $+1$ and flip $k$ leaves in a single branch on
the first side, where $k=1$ if $a=1$ and $1 \leq k \leq a-1$ if $a \geq 2$.

\begin{itemize}
\item When $a \geq 2$, the flipped branch contains both a $+1$ leaf and a $-1$
  leaf, so their difference vector detects $0$ (inner product $\pm 2$).
\item When $r \geq 2$, the flipped branch has signed leaf sum $a - 2k$ whereas
  every unflipped branch has leaf sum $a$. Choosing weights $t_1=1$ on the
  flipped branch and $t_2=-1$ on another branch on the same side gives inner
  product $-2k \neq 0$ at both $\pm\sqrt{a}$.
\item For every antisymmetric radial eigenvalue ($e=-1$), the unflipped sum
  vanishes by antisymmetry, so flipping $k$ leaves on the first side gives
  inner product $-2k \neq 0$.
\item For a symmetric radial eigenvalue $x$ ($e=+1$), the unflipped sum on each
  side is $F(x) = (x^2-a) + rx + ra$. Flipping $k$ leaves on the first side
  yields inner product $2(F(x) - k)$, where
  \[
    F(x) = x^2 - a + rx + ra.
  \]
\end{itemize}

When $a \geq 5$, the four integers $k \in \{1,2,3,4\}$ are all admissible
($1 \leq k \leq a-1$). Since $f_+$ has only three roots, at most three values
of $k$ can satisfy $F(x) = k$, leaving at least one admissible $k$ that detects
all three symmetric radial eigenvalues.

When $1 \leq a \leq 4$, take $k = 1$ and put
$q(x) = F(x) - 1 = x^2 - a + rx + ra - 1$. Both $f_+(x)$ and $q(x)$ are monic
in $x$ over $\mathbb{Z}[r,a]$, so their Sylvester resultant
$\Res_x(f_+, q) \in \mathbb{Z}[r,a]$ commutes with specialization of $(r,a)$.
Reducing modulo $r$ sets $r = 0$, giving
\[
  \Res_x(f_+, q)
  \equiv \Res_x\bigl((x-1)(x^2-a),\, x^2-a-1\bigr)
  = -a \pmod{r},
\]
because by multiplicativity of the resultant for monic polynomials the
specialized resultant is the product of $q|_{r=0}(1) = -a$ and the two values
$(x^2-a-1)|_{x^2=a} = -1$ at the roots of $x^2-a$. Consequently, $f_+$ and $q$
can share a root only when $r \mid a$. Table~\ref{tab:res-tra} lists
$\Res_x(f_+, q)$ for all eight pairs $(r,a)$ with $1 \leq a \leq 4$ and
$r \mid a$.

\begin{table}[ht]
\centering
\caption{Values of $\Res_x(f_+, q)$ for all divisor pairs $r \mid a$ with $1 \leq a \leq 4$.}\label{tab:res-tra}
\begin{tabular}{rrr}
\toprule
$a$ & $r$ & $\Res_x(f_+, q)$ \\
\midrule
1 & 1 & 1 \\
2 & 1 & 0 \\
2 & 2 & 16 \\
3 & 1 & 1 \\
3 & 3 & 465 \\
4 & 1 & 4 \\
4 & 2 & 274 \\
4 & 4 & 3256 \\
\bottomrule
\end{tabular}
\end{table}

Only $(r,a) = (1,2)$ has vanishing resultant. For $(r,a) = (1,2)$, flip one
leaf and its branch vertex on the first side (leaving all other signs $+1$).
The two differently signed leaves still detect $0$, and $r = 1$ has no branch
contrasts. At a symmetric radial eigenvalue $x$, subtracting $2(x+1)$ from
$2F(x)$ gives overlap $2(F(x) - x - 1) = 2(x^2 - 1)$; its potential zeros
$x = \pm 1$ are excluded by $f_+(1) = -1$ and $f_+(-1) = 3$. At an
antisymmetric radial eigenvalue $x$, the overlap is $-2(x+1)$, and its unique
zero $x = -1$ is excluded by $f_-(-1) = 1$. This completes the proof for every
$r,a \geq 1$.
\end{proof}

\section{The radial space of an arbitrary diameter-five tree}\label{sec:radial}

Every diameter-five tree has a unique central edge joining two centers
$L$ and $R$, each with at least one branch vertex of positive leaf count. For
each side $i \in \{L,R\}$, let $B_i \subset \mathbb{Z}_{\geq 1}$ be the
nonempty set of distinct branch sizes attached to center $i$, let
$b_{i,a} \geq 1$ be the number of branches of size $a \in B_i$, and let
$p_i \geq 0$ be the number of leaves adjacent directly to center $i$.

Consider the partition of the vertex set into the following cells:
\begin{itemize}
\item each center $L$ and $R$ as a singleton cell;
\item for each $i \in \{L,R\}$ and $a \in B_i$, one cell of all $b_{i,a}$
  branch vertices of size $a$ on side $i$, and a second cell of all
  $a b_{i,a}$ leaves attached to them;
\item for each $i \in \{L,R\}$ with $p_i > 0$, one cell of all $p_i$ direct
  leaves on side $i$.
\end{itemize}
This partition is equitable, so the subspace $W$ of vectors constant on each
cell is $A$-invariant and has dimension
\[
  D = 2 + 2(|B_L| + |B_R|) + \mathbf{1}_{p_L > 0} + \mathbf{1}_{p_R > 0}.
\]
Its orthogonal complement $W^\perp$ decomposes orthogonally into:
\begin{itemize}
\item leaf-difference spaces of eigenvalue $0$ (sum-zero vectors supported on
  the leaves of a single branch or on the direct leaves of one side), of total
  dimension
  $\sum_{i}\bigl((p_i - \mathbf{1}_{p_i>0}) + \sum_{a\in B_i} b_{i,a}(a-1)\bigr)$;
\item branch-contrast spaces of eigenvalues $\pm\sqrt{a}$ for each side $i$ and
  repeated size $a \in B_i$ with $b_{i,a} \geq 2$ (constructed with sum-zero
  branch weights $t_j$ as in Section~\ref{sec:spectrum-tra}), of total
  dimension $\sum_{i,a} 2(b_{i,a} - 1)$.
\end{itemize}
The dimensions of $W$, the leaf-difference spaces, and the branch-contrast
spaces sum to the tree order $2 + \sum_i p_i + \sum_{i,a} b_{i,a}(a+1)$.

\begin{lemma}\label{lem:radial-equations}
Let $v \in W$ be a nonzero eigenvector of $A$ with eigenvalue $x$, and let
$c_L, c_R$ denote its coordinates at centers $L,R$.
\begin{enumerate}
\item If $x \neq 0$, then:
  \begin{enumerate}
  \item on side $i$, the direct-leaf value (when $p_i > 0$) is $c_i / x$, and
    for every $a \in B_i$ with $x^2 \neq a$ the branch and branch-leaf values
    are $x c_i / (x^2-a)$ and $c_i / (x^2-a)$, respectively;
  \item if $x^2 = a \in B_i$, then $c_i = 0$;
  \item $(c_L, c_R) \neq (0,0)$;
  \item if $c_L = 0$, then $x^2 \in B_L \setminus B_R$ and $c_R \neq 0$ (and
    symmetrically with $L$ and $R$ interchanged).
  \end{enumerate}
\item If $x = 0$, then every branch value vanishes; $p_i > 0$ implies
  $c_i = 0$, and $p_i = 0$ implies $c_j = 0$ for the opposite center $j$. In
  particular, $0$ is not an eigenvalue of $A|_W$ whenever both sides have
  direct leaves ($p_L, p_R > 0$) or neither side does ($p_L = p_R = 0$).
\end{enumerate}
\end{lemma}

\begin{proof}
Let $y_{i,a}$, $z_{i,a}$, and $d_i$ (when $p_i > 0$) denote the common values
of $v$ on branch vertices of size $a$, their leaves, and the direct leaves on
side $i$. The eigenvalue equations $Av = xv$ on the cells are
\begin{align*}
  y_{i,a} &= x z_{i,a}, &
  c_i + a z_{i,a} &= x y_{i,a} \quad\text{(hence }(x^2-a)z_{i,a} = c_i\text{)}, \\
  c_i &= x d_i \quad (p_i > 0), &
  c_j + \sum_{a \in B_i} b_{i,a} y_{i,a} + p_i d_i &= x c_i \quad (\{i,j\} = \{L,R\}).
\end{align*}

Suppose $x \neq 0$. Parts~(a) and~(b) follow immediately from
$(x^2-a)z_{i,a} = c_i$, $y_{i,a} = x z_{i,a}$, and $c_i = x d_i$. For~(c), if
$c_L = c_R = 0$, then $d_i = 0$ and $y_{i,a} = z_{i,a} = 0$ for all
$a \neq x^2$; at most one group $a = x^2$ can remain on side $i$, and the
center-$i$ equation then forces $b_{i,x^2} y_{i,x^2} = 0$, hence
$y_{i,x^2} = z_{i,x^2} = 0$ and $v = 0$, a contradiction. For~(d), if
$c_L = 0$, then $c_R \neq 0$ by~(c), which forces $x^2 \notin B_R$ by~(b). If
$x^2 \notin B_L$, then all branch and direct-leaf values on side $L$ vanish by
(a), so the center-$L$ equation gives $c_R = 0$, contradicting~(c). Thus
$x^2 \in B_L \setminus B_R$.

Now suppose $x = 0$. The leaf equations give $y_{i,a} = 0$, the branch
equations give $z_{i,a} = -c_i/a$, the direct-leaf equation (when $p_i > 0$)
gives $c_i = 0$, and the center-$i$ equation reduces to $c_j + p_i d_i = 0$
(forcing $c_j = 0$ when $p_i = 0$). If either $p_L, p_R > 0$ or
$p_L = p_R = 0$, then $c_L = c_R = 0$, whence all $z_{i,a} = 0$ and all
$d_i = 0$, so $v = 0$.
\end{proof}

\subsection{Unequal branch sizes}\label{sec:unequal}

Set $m_i = \max(B_i \cup \{p_i\})$ for $i \in \{L,R\}$.

\begin{proposition}\label{prop:heterogeneous}
A diameter-five tree admits a full switching whenever
\[
  m_L \geq D + 2 \qquad\text{and}\qquad m_R \geq 2|B_L \setminus B_R| + 2,
\]
or whenever the same inequalities hold with $L$ and $R$ interchanged.
\end{proposition}

\begin{proof}
Choose one leaf group of size $m_i$ on each side $i \in \{L,R\}$ (either the
leaves of a single branch of size $m_i$ or the $p_i$ direct leaves at center
$i$). Start with all signs $+1$. For each repeated branch class ($b_{i,a} \geq 2$)
on either side, designate one representative branch; if the chosen leaf group
on side $i$ belongs to a repeated branch class, use that branch as its
representative. Flip one leaf of every representative other than those reserved
for the two chosen leaf groups. In each chosen leaf group, flip the first $k_i$
leaves in a fixed ordering, with $1 \leq k_i \leq m_i - 1$.

For every admissible pair $(k_L, k_R)$, the resulting sign vector $s$ detects
$0$ via a difference of two oppositely signed leaves in the chosen left group
(since $m_L \geq D + 2 \geq 8$), and detects every branch contrast
$\pm\sqrt{a}$ via the difference of signed leaf sums between the unique flipped
representative of class $(i,a)$ (which has $1$ or $k_i \geq 1$ flipped leaves)
and another, unflipped branch of the same class.

Now fix one eigenvector $v \in W$ for each distinct nonzero eigenvalue $x$ of
$A|_W$ (at most $D$ vectors). Because all leaves in the chosen group on side
$i$ lie in a single cell of $W$, $v$ takes a constant value $\zeta_i$ on that
group, and the inner product $s^T v$ has the form
$C_0 - 2\zeta_L k_L - 2\zeta_R k_R$ for a constant $C_0$ independent of
$(k_L, k_R)$.

Suppose $\zeta_L = 0$. By Lemma~\ref{lem:radial-equations}(1), this forces
$c_L = 0$ (if the chosen group is a branch of size $a_0 = x^2$, the center-$L$
equation gives $b_{L,a_0} x z_{L,a_0} = -c_R \neq 0$, so $\zeta_L \neq 0$ in
that case). Hence $x^2 \in B_L \setminus B_R$, $c_R \neq 0$, and
$x^2 \notin B_R$, which implies $\zeta_R \neq 0$. At most
$2|B_L \setminus B_R|$ eigenvalues $x$ satisfy $x^2 \in B_L \setminus B_R$. For
each such vector, $s^T v$ is independent of $k_L$ and affine in $k_R$ with
nonzero slope $-2\zeta_R$, forbidding at most one value of $k_R$. Since
$m_R - 1 \geq 2|B_L \setminus B_R| + 1$, we can fix an admissible $k_R$
avoiding all of these forbidden values.

For every remaining vector $v$ (at most $D$ in number), $\zeta_L \neq 0$, so
with $k_R$ fixed the overlap $s^T v$ is affine in $k_L$ with nonzero slope
$-2\zeta_L$, forbidding at most one value of $k_L$. Since
$m_L - 1 \geq D + 1 > D$, there exists an admissible $k_L$ avoiding all
forbidden values. The resulting sign vector $s$ detects every distinct
eigenvalue of $A$.
\end{proof}

\subsection{Reflected trees with unequal branch sizes}\label{sec:reflected-unequal}

Suppose the branch multiset and direct-leaf count agree on both sides, so that
$B_L = B_R$, $b_{L,a} = b_{R,a}$ for all $a \in B_L$, and $p_L = p_R$. Write
$q = |B_L|$, $p = p_L$, and $m = \max(B_L \cup \{p\})$.

\begin{proposition}\label{prop:reflected}
A reflected diameter-five tree admits a full switching whenever
\[
  m \geq 2q + 3 + \mathbf{1}_{p > 0}.
\]
\end{proposition}

\begin{proof}
The reflection involution swapping sides $L$ and $R$ splits $W$ into symmetric
and antisymmetric $A$-invariant subspaces $W = W_+ \oplus W_-$, each of
dimension $1 + 2q + \mathbf{1}_{p > 0}$. By Lemma~\ref{lem:radial-equations}(2),
neither parity contains the eigenvalue $0$. By
Lemma~\ref{lem:radial-equations}(1), no nonzero eigenvalue of $A|_{W_+}$ or
$A|_{W_-}$ can satisfy $c_L = 0$ or $x^2 \in B_L$ (since reflection symmetry
gives $c_R = \pm c_L$, so vanishing of one center forces both centers to
vanish). Consequently, every eigenvector in $W_+$ or $W_-$ has $c_L \neq 0$ and
takes a nonzero value on every leaf cell.

Prescribe center signs $+1$ on $L$ and $-1$ on $R$, and mirror all noncentral
signs across the two sides. For every eigenvector $v \in W_-$, the mirrored
noncentral contributions to $s^T v$ cancel in pairs, leaving
$s^T v = c_L - c_R = 2 c_L \neq 0$.

For the noncentral signs, choose one leaf group of size $m$ on side $L$ and its
mirror on side $R$, designate one representative per repeated branch class
(using the chosen group's branch if it lies in a repeated class), flip one leaf
in each unreserved representative on both sides, and flip $k$ leaves in the
chosen group on both sides with $1 \leq k \leq m - 1$. Since $m \geq 5$, the
resulting sign vector detects $0$ (by oppositely signed leaves in a chosen
group) and every branch contrast (by the signed leaf-sum differences). For
every eigenvector $v \in W_+$, the center signs cancel ($c_L - c_L = 0$) and
the common value $\zeta$ of $v$ on the chosen leaf group is nonzero, so
$s^T v$ is affine in $k$ with nonzero slope $-4\zeta$. Each of the at most
$\dim(W_+) = 1 + 2q + \mathbf{1}_{p>0}$ symmetric radial eigenvalues therefore
forbids at most one $k$. Because
$m - 1 \geq 2q + 2 + \mathbf{1}_{p > 0} = \dim(W_+) + 1$, at least one
admissible $k$ detects every symmetric radial eigenvalue.
\end{proof}

\section{All uniform reflected trees with direct leaves}\label{sec:direct}

For integers $r,a \geq 1$ and $p \geq 0$, let $T_r^p(a)$ denote the reflected
diameter-five tree obtained from $T_r(a)$ by attaching $p$ direct leaves to
each center.

\begin{theorem}\label{thm:uniform-direct}
For all integers $r,a \geq 1$ and $p \geq 0$, the tree $T_r^p(a)$ admits a sign
switching that makes every distinct adjacency eigenvalue main.
\end{theorem}

\begin{proof}
Theorem~\ref{thm:uniform} handles $p = 0$. For $p > 0$,
Proposition~\ref{prop:reflected} with $q = 1$ handles every case with
$\max(a,p) \geq 6$. It remains to treat the $25$ pairs with
$1 \leq a,p \leq 5$ and arbitrary $r \geq 1$.

Set the center signs to $+1$ and $-1$, flip one branch leaf on each side, and
flip one direct leaf on each side if $a = 1$ and $p \geq 2$; all noncentral
signs are mirrored across the two sides. By the argument of
Proposition~\ref{prop:reflected}, this sign vector detects every antisymmetric
radial eigenvalue (via the center signs) and every branch contrast when
$r \geq 2$ (via the single altered branch on each side). For eigenvalue $0$:
$0$ is absent from $W$ by Lemma~\ref{lem:radial-equations}(2), and if $a=p=1$
there are no leaf-difference modes at all; otherwise either $a \geq 2$ (so the
flipped branch has oppositely signed leaves) or $a=1, p \geq 2$ (so the direct
leaves have opposite signs), detecting $0$.

It remains to verify the symmetric radial eigenvalues in $W_+$. Normalize an
eigenvector $v \in W_+$ to have $c_L = c_R = 1$. By
Lemma~\ref{lem:radial-equations}(1), $x \neq 0$ and $x^2 \neq a$, and the
center equation $x - 1 = rx/(x^2-a) + p/x$ clears to
\begin{equation}\label{eq:g}
  g(x) = (x - 1)x(x^2 - a) - r x^2 - p(x^2 - a) = 0.
\end{equation}
The center signs cancel in $s^T v$, so $\frac{1}{2} s^T v$ equals the signed
noncentral sum on side $L$. Using the center equation
$rx/(x^2-a) = x - 1 - p/x$ and $ra/(x^2-a) = x \cdot rx/(x^2-a) - r$, the
unflipped noncentral sum on side $L$ simplifies to
\[
  \frac{r(x+a)}{x^2-a} + \frac{p}{x}
  = (x - 1) + x\left(x - 1 - \frac{p}{x}\right) - r
  = x^2 - (r + p + 1).
\]
Set $\ell = 1$ if $a = 1$ and $p \geq 2$, and $\ell = 0$ otherwise. Subtracting
$2/(x^2-a) + 2\ell/x$ for the flipped leaf (or leaves) and multiplying
$\frac{1}{2} s^T v$ by $x(x^2-a)$ gives
\begin{equation}\label{eq:h}
  h(x) = \bigl(x^2 - (r + p + 1)\bigr)x(x^2 - a) - 2x - 2\ell(x^2 - a).
\end{equation}
The factor $x(x^2-a)$ never vanishes at a root of $g$, since $g(0) = ap > 0$
and $g(\pm\sqrt{a}) = -ra < 0$.

For each fixed pair $(a,p) \in \{1,\dots,5\}^2$, both $g$ and $h$ are monic in
$x$ over $\mathbb{Z}[r]$, so their resultant
$R_{a,p}(r) = \Res_x(g,h) \in \mathbb{Z}[r]$ commutes with specialization of
$r$. Thus $g$ and $h$ share a root at a positive integer $r$ if and only if
$R_{a,p}(r) = 0$. From the primitive part of $R_{a,p}(r)$, remove the linear
factor $r - 1$ at $(a,p) = (2,1)$ and the factor $r^2$ at $(a,p) = (2,2)$,
where the unreduced resultants factor as
\[
  R_{2,1}(r) = -8(r - 1)(r^4 + 11r^2 - 1), \qquad
  R_{2,2}(r) = -16 r^2(r^3 + 2r^2 + 14r - 18).
\]
For every $(a,p) \in \{1,\dots,5\}^2$, the remaining primitive integer
polynomial is nonzero at every residue modulo the integer $v$ listed in
Table~\ref{tab:mod-certs}, hence has no integer roots at all. (The accompanying
Lean~4 module \path{formalization/ModularExclusions.lean} formally verifies
all $25$ modular nonvanishing certificates.) Therefore $R_{a,p}(r)$ has a
positive integer root only for $(r,a,p) = (1,2,1)$.

\begin{table}[ht]
\centering
\caption{Moduli $v$ excluding integer roots of the reduced primitive resultants for $1 \leq a,p \leq 5$.}\label{tab:mod-certs}
\begin{tabular}{rccccc}
\toprule
$a \setminus p$ & 1 & 2 & 3 & 4 & 5 \\
\midrule
1 & 3 & 3 & 3 & 5 & 3 \\
2 & 2 & 5 & 7 & 7 & 17 \\
3 & 3 & 3 & 3 & 3 & 3 \\
4 & 3 & 3 & 5 & 3 & 3 \\
5 & 5 & 5 & 5 & 5 & 5 \\
\bottomrule
\end{tabular}
\end{table}

For the sole exception $(r,a,p) = (1,2,1)$, keep both center signs equal to
$+1$ and flip one branch leaf and the direct leaf on the same side. There are
no repeated-branch contrasts ($r=1$), and the two differently signed leaves of
the flipped branch detect $0$. For a symmetric radial eigenvalue ($c_L=c_R=1$),
equation~\eqref{eq:g} becomes $g(x) = (x-1)(x^3 - 4x - 2)$. The unflipped total
sum on both sides is $2\bigl(1 + (x^2 - 3)\bigr) = 2(x^2 - 2)$, so flipping one
branch leaf and one direct leaf on one side gives half-overlap multiplied by
$x(x^2-2)$ equal to
\[
  H(x) = x(x^2 - 2)^2 - x - (x^2 - 2).
\]
We have $H(1) = 1 \neq 0$, and reducing $H(x)$ modulo $x^3 - 4x - 2$ yields
$(x+1)(x+2)$, while $x^3 - 4x - 2$ takes values $1$ and $-2$ at $x = -1$ and
$x = -2$, respectively. For an antisymmetric radial eigenvalue
($c_L = 1, c_R = -1$), the unflipped sum vanishes, so the overlap is
$-2\bigl(\frac{1}{x^2-2} + \frac{1}{x}\bigr) = -2\frac{(x-1)(x+2)}{x(x^2-2)}$.
Its potential zeros $x = 1, -2$ are excluded because the antisymmetric radial
polynomial $(x+1)x(x^2-2) - x^2 - (x^2-2)$ evaluates to $-2$ at both $1$ and
$-2$. This completes the proof.
\end{proof}

\section{Exact constraints for arbitrary diameter-five trees}\label{sec:remaining}

For an arbitrary diameter-five tree with parameters $(B_i, b_{i,a}, p_i)$, a
sign vector is summarized on $W$ and its orthogonal complement by center signs
$\sigma_i \in \{-1,1\}$, branch signs $\beta_{i,a,j} \in \{-1,1\}$
($1 \leq j \leq b_{i,a}$), branch leaf-sign sums
$u_{i,a,j} \in \{-a, -a+2, \dots, a\}$, and direct-leaf sign sums
$v_i \in \{-p_i, -p_i+2, \dots, p_i\}$. Every tuple of such integers is
independently realizable by a sign vector $s$. Define
\[
  U_{i,a}(x) = \sum_{j=1}^{b_{i,a}} (x\beta_{i,a,j} + u_{i,a,j}), \quad
  H_i(x) = x - \frac{p_i}{x} - \sum_{a\in B_i} \frac{b_{i,a} x}{x^2 - a}, \quad
  L_i(x) = \sigma_i + \frac{v_i}{x} + \sum_{a\in B_i} \frac{U_{i,a}(x)}{x^2 - a}.
\]

\begin{proposition}\label{prop:constraints}
A sign vector with parameters $(\sigma_i, \beta_{i,a,j}, u_{i,a,j}, v_i)$ is a
full switching for a diameter-five tree if and only if the following conditions
hold:
\begin{enumerate}
\item For every $x \neq 0$ with $x^2 \notin B_L \cup B_R$ and
  $H_L(x) H_R(x) = 1$, we have $L_L(x) + H_L(x) L_R(x) \neq 0$.
\item For every $x \neq 0$ with $x^2 = a \in B_L \setminus B_R$ and
  $H_R(x) = 0$, either $L_R(x) - U_{L,a}(x)/(b_{L,a} x) \neq 0$ or
  $x\beta_{L,a,j} + u_{L,a,j}$ varies with $j$ (and symmetrically for
  $a \in B_R \setminus B_L$).
\item For every $a \in B_i$ with $b_{i,a} \geq 2$ and $x = \pm\sqrt{a}$, either
  $x\beta_{i,a,j} + u_{i,a,j}$ varies with $j$ on at least one side $i$ with
  $b_{i,a} \geq 2$, or $x$ is detected by the pole radial eigenvector of~(2)
  when $a$ lies exclusively on one side and the opposite $H$-function vanishes
  at $x$.
\item If $0$ is an eigenvalue of $A$ (which occurs when some leaf group has
  size at least $2$ or when exactly one side has direct leaves), then either
  $|u_{i,a,j}| < a$ or $|v_i| < p_i$ for some group, or else (when exactly one
  side, say $L$, has direct leaves)
  $\sigma_R - v_L/p_L - \sum_{a\in B_R}\sum_{j=1}^{b_{R,a}} u_{R,a,j}/a \neq 0$
  (and symmetrically when $p_R > 0, p_L = 0$).
\end{enumerate}
\end{proposition}

\begin{proof}
By the orthogonal decomposition of Section~\ref{sec:radial}, each distinct
eigenvalue $x$ is detected if and only if $s$ has nonzero inner product with at
least one eigenvector of $A$ at $x$ lying in $W$, a branch-contrast space, or a
leaf-difference space.

For $x \neq 0$ with $x^2 \notin B_L \cup B_R$, the center equations in
Lemma~\ref{lem:radial-equations}(1) read $c_R = H_L(x) c_L$ and
$c_L = H_R(x) c_R$, giving a unique radial direction $(c_L, c_R) = (1, H_L(x))$
precisely when $H_L(x) H_R(x) = 1$, with overlap $L_L(x) + H_L(x) L_R(x)$.

For $x \neq 0$ at a pole $x^2 = a \in B_L \setminus B_R$,
Lemma~\ref{lem:radial-equations}(1) forces $(c_L, c_R) = (0,1)$ (up to scale),
$H_R(x) = 0$, and left branch-leaf value $z_{L,a} = -1/(b_{L,a} x)$ in group
$a$, yielding radial overlap $L_R(x) - U_{L,a}(x)/(b_{L,a} x)$. At a shared
pole $x^2 \in B_L \cap B_R$, Lemma~\ref{lem:radial-equations}(1) excludes
radial eigenvectors. On any side $i$ with $b_{i,a} \geq 2$, a branch-contrast
vector with weights $\sum_j t_j = 0$ has inner product
$\sum_j t_j(x\beta_{i,a,j} + u_{i,a,j})$, which can be made nonzero if and only
if $x\beta_{i,a,j} + u_{i,a,j}$ is not constant in $j$.

At $x = 0$, a leaf-difference eigenvector detects $0$ if and only if some leaf
group contains both $+1$ and $-1$ signs, i.e., $|u_{i,a,j}| < a$ or
$|v_i| < p_i$. By Lemma~\ref{lem:radial-equations}(2), a radial $0$-eigenvector
exists precisely when exactly one side has direct leaves; if $p_L > 0$ and
$p_R = 0$, normalizing $(c_L, c_R) = (0,1)$ yields $d_L = -1/p_L$,
$z_{R,a} = -1/a$, and all other coordinates $0$, giving the stated overlap.
\end{proof}

The general diameter-five conjecture remains open outside the families covered
by Theorems~\ref{thm:uniform} and~\ref{thm:uniform-direct} and
Propositions~\ref{prop:heterogeneous} and~\ref{prop:reflected}---for example,
reflected trees with several distinct small branch sizes and nonreflected trees
with small leaf groups.

\section{Computational verification and formalization}\label{sec:repro}

The companion repository provides exact verification scripts
(\path{replay.py}, using Python~3.12 and SymPy~1.14.0) and Lean~4.34.1
formalization modules (\path{formalization/FlipChoices.lean} and
\path{formalization/ModularExclusions.lean}):

\begin{itemize}
\item Exact Krylov-rank checks against the square-free characteristic
  polynomial degree verify the small-graph switching counts ($K_1, K_2, P_3,
  P_4, K_4-e, P_6, K_{1,3}, C_6, K_{3,3}$, and the triangular prism), which are
  independently cross-checked via rational primary spectral subspaces in
  \path{verification/switching_controls.py}.
\item Exact symbolic arithmetic verifies the resultant congruence and divisor
  table of Theorem~\ref{thm:uniform} and tests the constructed sign vectors
  directly on all $40$ adjacency matrices $T_r(a)$ with $1 \leq r \leq 5$ and
  $1 \leq a \leq 8$.
\item Exact resultant computation over $\mathbb{Z}[r]$ confirms that
  $(r,a,p) = (1,2,1)$ is the unique positive integer root across all $25$ pairs
  $1 \leq a,p \leq 5$ in Theorem~\ref{thm:uniform-direct}, verifies the
  exceptional witness on the adjacency matrix of $T_1^1(2)$, and finds exact
  witnesses for $180$ heterogeneous diameter-five parameter rows with one or
  two branches per side, branch sizes in $\{1,2,3\}$, and $p_i \in \{0,1\}$.
\item Lean~4.34.1 verifies the four-choice counting step
  (\texttt{Switching.four\_choices}) and all $25$ modular nonvanishing
  certificates (\texttt{Switching.reflected\_exclusions}). The script
  \texttt{verification/\allowbreak check\_lean\_inputs.py} cross-checks the $25$ polynomial
  coefficient lists in Lean against the recomputed resultants.
\end{itemize}

\begin{thebibliography}{9}

\bibitem{AKMP}
S.~Akbari, H.~Kumar, B.~Mohar, and S.~Pragada,
arXiv:2609.27046, Section~1.4 (Problem~1.6),
\url{https://arxiv.org/html/2609.27046#S1.SS4}.

\bibitem{ChaShan}
H.~Cha and H.~Shan,
arXiv:2607.13577v1, Section~3 (diameter-five tree notation and families $T_r(a)$, $T_r^p(a)$),
\url{https://arxiv.org/html/2607.13577v1#S3}.

\bibitem{Dong}
H.~Dong,
Paper~8 (diameter-at-most-four switching and Krylov certificate),
\url{https://github.com/Horace-Maxwell/ai4math-results/blob/main/papers/switching-diam4/note.tex};
Lean statements:
\url{https://raw.githubusercontent.com/Horace-Maxwell/ai4math-results/main/lean/Research/Backfill/Paper8/Challenge.lean}.

\end{thebibliography}

\end{document}
